\documentclass{article}

% typesetting and text macros
\usepackage[letterpaper, textheight=9.0in, textwidth=5.2in]{geometry}
\usepackage{setspace}
\setstretch{1.08}
\usepackage{fancyhdr}
\setlength{\topmargin}{0.0in}
\setlength{\headheight}{3.0ex}\addtolength{\topmargin}{-0.5\headheight} % make room for headers
\setlength{\headsep}{3.0ex}\addtolength{\topmargin}{-0.5\headsep}
\pagestyle{fancy}
%\renewcommand{\headrulewidth}{0.0pt}
\fancyhf{}
\fancyhead[R]{\sffamily\thepage}
\sloppy\sloppypar\raggedbottom\frenchspacing

\usepackage{amsmath}
\DeclareMathOperator{\ee}{e}
\DeclareMathOperator{\ii}{i}
\DeclareMathOperator*{\argmin}{argmin}
\DeclareMathOperator{\loss}{loss}
\DeclareMathOperator{\rank}{rank}

\fancyhead[L]{\sffamily team / fully continuous heteroskedastic data factorizations for arbitrarily ragged data}
\title{\bfseries A continuous spatial map of coefficients, controlling a learned set of continuous eigenfunctions, as a model for heterogeneous and ragged spectral imaging data, by way of robust heteroskedastic factorizations in function space, or, a fully functional version of robust heteroskedastic matrix factorization}
\author{Hogg \and others}
\date{August 2026}

\begin{document}

\maketitle

\paragraph{Abstract:}
The original dimensionality reduction method is PCA, which represents a rectangular data matrix with a rank-$K$ approximation, optimizing a mean-squared loss.
There are extensions and improvements, but most of them require rectangular data (though some permit missing values).
In many astrophysical contexts, data are truly ragged; for example,
think of \textsl{LSST} or NASA \textsl{SPHEREx} observations:
Each small sky region gets a different number of observations, at different locations both spatially (in celestial coordinates) and in time and wavelength.
Here we provide a dimensionality reduction that learns a basis of $K$ continuous functions of wavelength, outer producted with $K$ continuous functions of sky position that delivers an accurate, technically robust, low-dimensional representation of the ragged data.
The method interpolates and extrapolates spectra, answers counterfactual questions about missing data, and identifies (and down-weights) outlier and anomalous data points.
This discussion is aimed at \textsl{SPHEREx} data analysis, but generalizations could be used in a wide range of scientific and engineering applications.

\paragraph{Introduction:}
First there was principal components analysis (PCA) \cite{pca},
then there was heteroskedastic matrix factorization (HMF) \cite{hmf},
and then there was Robust HMF \cite{rhmf}.
In all of these, the idea is that the user has \emph{rectangular data}:
There are $N$ data vectors (lists really) $\vec{y}_n$, each of which is an ordered list of $d$ features, making up a $N\times d$ data matrix $Y$.
We will refer to a single element of $Y$ as $y_{nj}$ with $1\leq n\leq N$ and $1\leq j\leq d$.

In PCA, each of these $N\times d$ data points $y_{nj}$ is treated as identically valuable;
in HMF, each data point comes with its own weight $w_{nj}$ (and, indeed, some of those weights can be zero);
in robust HMF, each data point gets a weight adjustment depending on how well the low-rank structure that I am about to describe fits that data point.
Importantly, HMF and robust HMF permit there to be missing data, because data points can be given zero weights.

In all three cases, the fundamental model is a low-rank linear representation, or a matrix factorization.
That is, each model can be written as
\begin{align}
    y_{nj} &= \sum_{k=1}^K a_{nk}\,g_{kj} + s_{nj} ~,
\end{align}
where the $a_{nk}$ are coefficients ($K$ for each data vector $n$,
the $g_{kj}$ are elements of eigenvectors or eigenlists, which can be thought of as comprising $K$ $d$-vectors $\vec{g}_k$,
and the $s_{nj}$ are per-measurement residuals.
Some might write this equation as
\begin{align}
    Y &= A\,G + S~,
\end{align}
where $A$, $G$, and $S$ are $N\times K$, $K\times d$, and $N\times d$ matrices of coefficients, eigenvectors, and residuals.
The only differences, in this context, between PCA, HMF, and robust HMF, is what loss function of the residuals is used to set the coefficients $A$ and eigenvectors $G$.
That is, each of these methods can be written in the form
\begin{align}
    \hat{A}, \hat{G} &= \argmin_{A, G} \loss(Y - A\,G) \quad\text{s. t.}~\rank(A\,G) = K ~,
\end{align}
for some loss (different for each method).

In many astronomical contexts, we don't get rectangular data.
For example, in \textsl{LSST} \cite{lsst}, each region on the sky will have a set of observations taken at different detailed sky positions, at different times $t_{nj}$ and through different filters.
In NASA \textsl{SPHEREx} \cite{spherex}, no two regions on the sky will be observed at the same list of wavelengths, nor at the same detailed sky positions within the region; indeed different regions get different numbers of observations, and each observation $j$ is at a different sky position and at a different value of (effective) wavelength (in general).
Can we nonetheless produce a dimensionally reduced, continuous map of the specific intensity?
That is, can we benefit from the dimensionality reduction and de-noising of PCA and HMF and RHMF without requiring data that are rectangular or similarly structured?
For specificity here, we will proceed in the \textsl{SPHEREx} context.

\paragraph{Setup:}
Now the setup is that we have $N$ measurements $y_j$ of the sky, each of which represents a single-pixel, calibrated measurement of specific intensity from one of the pixels in one of the SPHEREx cameras.
We are thinking of the imaging data as a huge, unstructured set of pixel measurements.
Each meassurement $y_j$ has associated with it a two-dimensional (effective, or central) sky position $x_j$, a one-dimensional (effective, or central) wavelength $\lambda_j$, a statistical weight or inverse noise variance estimate $w_j$, and lots of other housekeeping data, including most notably for what follows, the (effective, or central) time $t_j$ at which the measurement was made, in some sensible time system (barycentric, perhaps).

The idea is that we are going to learn a continuous, rank-$K$ model for these data, such that the data are well represented or explained as evaluations of this low-rank functional model.
In full generality, we seek $K$ functions $a_k(x)$ and $K$ functions $g_k(\lambda)$
\begin{align}
    y_j &= \sum_{k=1}^K a_k(x_j)\,g_k(\lambda_j) + s_j
\end{align}
such that a loss function acting on the residuals $s_j$ and weights $w_j$ is minimized, perhaps also subject to additional constraints such as orthogonality or non-negativity.

Now how to learn continuous functions?
Choose a basis and learn basis coefficients, probably.
So this will look something like
\begin{align}
    a_k(x) &= \sum_\ell \alpha_{k\ell}\,B_\ell(x) \\
    g_k(\lambda) &= \sum_m \gamma_{km}\,H_m(\lambda) ~,
\end{align}
where the $B_\ell()$ comprise a sensible, generic basis of spatial functions,
and the $H_m()$ comprise a sensible basis of wavelength functions.
Once we choose our bases, our loss function, and our constraints, all we need do is optimize.

\paragraph{Choice of bases:}

\paragraph{Loss functions for robustness:}

\paragraph{Method:}

\bibliographystyle{plain}
\bibliography{ragged}
\end{document}
